\section{Going up and going down}
\label{section-going-up}

\noindent
Suppose $\mathfrak p$, $\mathfrak p'$ are primes
of the ring $R$. Let $X = \Spec(R)$ with the Zariski
topology. Denote $x \in X$ the point corresponding
to $\mathfrak p$ and $x' \in X$ the point corresponding
to $\mathfrak p'$. Then we have:
$$
x' \leadsto x \Leftrightarrow \mathfrak p' \subset \mathfrak p.
$$
In words: $x$ is a specialization of $x'$ if and
only if $\mathfrak p' \subset \mathfrak p$.
See Topology, Section \ref{topology-section-specialization}
for terminology and notation.

\begin{definition}
\label{definition-going-up-down}
Let $\varphi : R \to S$ be a ring map.
\begin{enumerate}
\item We say a $\varphi : R \to S$ satisfies {\it going up} if
given primes $\mathfrak p \subset \mathfrak p'$ in $R$
and a prime $\mathfrak q$ in $S$ lying over $\mathfrak p$
there exists a prime $\mathfrak q'$ of $S$ such that
(a) $\mathfrak q \subset \mathfrak q'$, and (b)
$\mathfrak q'$ lies over $\mathfrak p'$.
\item We say a $\varphi : R \to S$ satisfies {\it going down} if
given primes $\mathfrak p \subset \mathfrak p'$ in $R$
and a prime $\mathfrak q'$ in $S$ lying over $\mathfrak p'$
there exists a prime $\mathfrak q$ of $S$ such that
(a) $\mathfrak q \subset \mathfrak q'$, and (b)
$\mathfrak q$ lies over $\mathfrak p$.
\end{enumerate}
\end{definition}

\noindent
So far we have see the following cases of this:
\begin{enumerate}
\item An integral ring map satisfies going up, see
Lemma \ref{lemma-integral-going-up}.
\item As a special case finite ring maps satisfy going up.
\item As a special case quotient maps $R \to R/I$ satisfy going up.
\item A flat ring map satisfies going down, see
Lemma \ref{lemma-flat-going-down}
\item As a special case any localization satisfies going down.
\item An extension $R \subset S$ of domains, with $R$ normal
and $S$ integral over $R$ satisfies going down, see
Proposition \ref{proposition-going-down-normal-integral}.
\end{enumerate}
Here is another case where going down holds.

\begin{lemma}
\label{lemma-open-going-down}
Let $R \to S$ be a ring map. If the induced map
$\varphi : \Spec(S) \to \Spec(R)$ is open, then
$R \to S$ satisfies going down.
\end{lemma}

\begin{proof}
Suppose that $\mathfrak p \subset \mathfrak p' \subset R$ and
$\mathfrak q' \subset S$ lies over $\mathfrak p'$. As $\varphi$ is open,
for every $g \in S$, $g \not \in \mathfrak q'$ we see that $\mathfrak p$
is in the image of $D(g) \subset \Spec(S)$. In other words
$S_g \otimes_R \kappa(\mathfrak p)$ is not zero. Since $S_{\mathfrak q'}$
is the directed colimit of these $S_g$ this implies
that $S_{\mathfrak q'} \otimes_R \kappa(\mathfrak p)$ is not
zero, see
Lemmas \ref{lemma-localization-colimit} and
\ref{lemma-tensor-products-commute-with-limits}.
Hence $\mathfrak p$ is in the image of
$\Spec(S_{\mathfrak q'}) \to \Spec(R)$ as desired.
\end{proof}

\begin{lemma}
\label{lemma-going-up-down-specialization}
Let $R \to S$ be a ring map.
\begin{enumerate}
\item $R \to S$ satisfies going down if and only if
generalizations lift along the map $\Spec(S) \to \Spec(R)$,
see Topology, Definition \ref{topology-definition-lift-specializations}.
\item $R \to S$ satisfies going up if and only if
specializations lift along the map $\Spec(S) \to \Spec(R)$,
see Topology, Definition \ref{topology-definition-lift-specializations}.
\end{enumerate}
\end{lemma}

\begin{proof}
Omitted.
\end{proof}

\begin{lemma}
\label{lemma-going-up-down-composition}
Suppose $R \to S$ and $S \to T$ are ring maps satisfying
going down. Then so does $R \to T$. Similarly for going up.
\end{lemma}

\begin{proof}
According to Lemma \ref{lemma-going-up-down-specialization}
this follows from
Topology, Lemma \ref{topology-lemma-lift-specialization-composition}
\end{proof}

\begin{lemma}
\label{lemma-image-stable-specialization-closed}
Let $R \to S$ be a ring map. Let $T \subset \Spec(R)$
be the image of $\Spec(S)$. If $T$ is stable under specialization,
then $T$ is closed.
\end{lemma}

\begin{proof}
We give two proofs.

\medskip\noindent
First proof. Let $\mathfrak p \subset R$ be a prime ideal such that
the corresponding point of $\Spec(R)$ is in the closure
of $T$. This means that for every $f \in R$, $f \not \in \mathfrak p$
we have $D(f) \cap T \not = \emptyset$. Note that $D(f) \cap T$
is the image of $\Spec(S_f)$ in $\Spec(R)$. Hence
we conclude that $S_f \not = 0$. In other words, $1 \not = 0$ in
the ring $S_f$. Since $S_{\mathfrak p}$ is the directed colimit
of the rings $S_f$ we conclude that $1 \not = 0$ in
$S_{\mathfrak p}$. In other words, $S_{\mathfrak p} \not = 0$ and
considering the image of $\Spec(S_{\mathfrak p})
\to \Spec(S) \to \Spec(R)$ we see there exists
a $\mathfrak p' \in T$ with $\mathfrak p' \subset \mathfrak p$.
As we assumed $T$ closed under specialization we conclude $\mathfrak p$
is a point of $T$ as desired.

\medskip\noindent
Second proof. Let $I = \Ker(R \to S)$. We may replace $R$ by $R/I$.
In this case the ring map $R \to S$ is injective.
By Lemma \ref{lemma-injective-minimal-primes-in-image}
all the minimal primes of $R$ are contained in the image $T$. Hence
if $T$ is stable under specialization then it contains all primes.
\end{proof}

\begin{lemma}
\label{lemma-going-up-closed}
Let $R \to S$ be a ring map. The following are equivalent:
\begin{enumerate}
\item Going up holds for $R \to S$, and
\item the map $\Spec(S) \to \Spec(R)$ is closed.
\end{enumerate}
\end{lemma}

\begin{proof}
It is a general fact that specializations lift along a
closed map of topological spaces, see
Topology, Lemma \ref{topology-lemma-closed-open-map-specialization}.
Hence the second condition implies the first.

\medskip\noindent
Assume that going up holds for $R \to S$.
Let $V(I) \subset \Spec(S)$ be a closed set.
We want to show that the image of $V(I)$ in $\Spec(R)$ is closed.
The ring map $S \to S/I$ obviously satisfies going up.
Hence $R \to S \to S/I$ satisfies going up,
by Lemma \ref{lemma-going-up-down-composition}.
Replacing $S$ by $S/I$ it suffices to show the image $T$
of $\Spec(S)$ in $\Spec(R)$ is closed.
By Topology, Lemmas \ref{topology-lemma-open-closed-specialization}
and \ref{topology-lemma-lift-specializations-images} this
image is stable under specialization. Thus the result follows
from Lemma \ref{lemma-image-stable-specialization-closed}.
\end{proof}

\begin{lemma}
\label{lemma-constructible-stable-specialization-closed}
Let $R$ be a ring. Let $E \subset \Spec(R)$ be a constructible subset.
\begin{enumerate}
\item If $E$ is stable under specialization, then $E$ is closed.
\item If $E$ is stable under generalization, then $E$ is open.
\end{enumerate}
\end{lemma}

\begin{proof}
First proof. The first assertion
follows from Lemma \ref{lemma-image-stable-specialization-closed}
combined with Lemma \ref{lemma-constructible-is-image}.
The second follows because the complement of a constructible
set is constructible
(see Topology, Lemma \ref{topology-lemma-constructible}),
the first part of the lemma and Topology,
Lemma \ref{topology-lemma-open-closed-specialization}.

\medskip\noindent
Second proof. Since $\Spec(R)$ is a spectral space by
Lemma \ref{lemma-spec-spectral} this is a special case of
Topology, Lemma
\ref{topology-lemma-constructible-stable-specialization-closed}.
\end{proof}

\begin{proposition}
\label{proposition-fppf-open}
Let $R \to S$ be flat and of finite presentation.
Then $\Spec(S) \to \Spec(R)$ is open.
More generally this holds for any ring map $R \to S$ of
finite presentation which satisfies going down.
\end{proposition}

\begin{proof}
If $R \to S$ is flat, then $R \to S$ satisfies going down by
Lemma \ref{lemma-flat-going-down}. Thus to prove the lemma we may
assume that $R \to S$ has finite presentation and satisfies going down.

\medskip\noindent
Since the standard opens $D(g) \subset \Spec(S)$, $g \in S$
form a basis for the topology, it suffices to prove that the
image of $D(g)$ is open. Recall that $\Spec(S_g) \to \Spec(S)$
is a homeomorphism of $\Spec(S_g)$ onto $D(g)$
(Lemma \ref{lemma-standard-open}).
Since $S \to S_g$ satisfies going down (see above), we see that
$R \to S_g$ satisfies going down by
Lemma \ref{lemma-going-up-down-composition}. Thus after replacing
$S$ by $S_g$ we see it suffices to prove the image is open.
By Chevalley's theorem (Theorem \ref{theorem-chevalley})
the image is a constructible set $E$. And $E$ is stable
under generalization because $R \to S$ satisfies going down,
see Topology, Lemmas \ref{topology-lemma-open-closed-specialization}
and \ref{topology-lemma-lift-specializations-images}.
Hence $E$ is open by
Lemma \ref{lemma-constructible-stable-specialization-closed}.
\end{proof}

\begin{lemma}
\label{lemma-same-image}
Let $k$ be a field, and let $R$, $S$ be $k$-algebras.
Let $S' \subset S$ be a sub $k$-algebra, and let $f \in S' \otimes_k R$.
In the commutative diagram
$$
\xymatrix{
\Spec((S \otimes_k R)_f) \ar[rd] \ar[rr] & &
\Spec((S' \otimes_k R)_f) \ar[ld] \\
& \Spec(R) &
}
$$
the images of the diagonal arrows are the same.
\end{lemma}

\begin{proof}
Let $\mathfrak p \subset R$ be in the image of the south-west
arrow. This means (Lemma \ref{lemma-in-image}) that
$$
(S' \otimes_k R)_f \otimes_R \kappa(\mathfrak p)
=
(S' \otimes_k \kappa(\mathfrak p))_f
$$
is not the zero ring, i.e., $S' \otimes_k \kappa(\mathfrak p)$
is not the zero ring and the image of $f$ in it is not nilpotent.
The ring map
$S' \otimes_k \kappa(\mathfrak p) \to S \otimes_k \kappa(\mathfrak p)$
is injective. Hence also $S \otimes_k \kappa(\mathfrak p)$
is not the zero ring and the image of $f$ in it is not nilpotent.
Hence $(S \otimes_k R)_f \otimes_R \kappa(\mathfrak p)$
is not the zero ring. Thus (Lemma \ref{lemma-in-image})
we see that $\mathfrak p$ is in the image of the south-east arrow
as desired.
\end{proof}

\begin{lemma}
\label{lemma-map-into-tensor-algebra-open}
Let $k$ be a field.
Let $R$ and $S$ be $k$-algebras.
The map $\Spec(S \otimes_k R) \to \Spec(R)$
is open.
\end{lemma}

\begin{proof}
Let $f \in S \otimes_k R$.
It suffices to prove that the image of the standard open $D(f)$ is open.
Let $S' \subset S$ be a finite type $k$-subalgebra such that
$f \in S' \otimes_k R$. The map $R \to S' \otimes_k R$ is flat
and of finite presentation, hence the image $U$ of
$\Spec((S' \otimes_k R)_f) \to \Spec(R)$ is open
by Proposition \ref{proposition-fppf-open}.
By Lemma \ref{lemma-same-image} this is also the image of $D(f)$ and we win.
\end{proof}

\noindent
Here is a tricky lemma that is sometimes useful.

\begin{lemma}
\label{lemma-unique-prime-over-localize-below}
Let $R \to S$ be a ring map.
Let $\mathfrak p \subset R$ be a prime.
Assume that
\begin{enumerate}
\item there exists a unique prime $\mathfrak q \subset S$ lying over
$\mathfrak p$, and
\item either
\begin{enumerate}
\item going up holds for $R \to S$, or
\item going down holds for $R \to S$ and there is at most one prime
of $S$ above every prime of $R$.
\end{enumerate}
\end{enumerate}
Then $S_{\mathfrak p} = S_{\mathfrak q}$.
\end{lemma}

\begin{proof}
Consider any prime $\mathfrak q' \subset S$ which corresponds to
a point of $\Spec(S_{\mathfrak p})$. This means that
$\mathfrak p' = R \cap \mathfrak q'$ is contained in $\mathfrak p$.
Here is a picture
$$
\xymatrix{
\mathfrak q' \ar@{-}[d] \ar@{-}[r] & ? \ar@{-}[r] \ar@{-}[d] & S \ar@{-}[d] \\
\mathfrak p' \ar@{-}[r] & \mathfrak p \ar@{-}[r] & R
}
$$
Assume (1) and (2)(a).
By going up there exists a prime $\mathfrak q'' \subset S$
with $\mathfrak q' \subset \mathfrak q''$ and $\mathfrak q''$
lying over $\mathfrak p$. By the uniqueness of $\mathfrak q$ we
conclude that $\mathfrak q'' = \mathfrak q$. In other words
$\mathfrak q'$ defines a point of $\Spec(S_{\mathfrak q})$.

\medskip\noindent
Assume (1) and (2)(b).
By going down there exists a prime $\mathfrak q'' \subset \mathfrak q$
lying over $\mathfrak p'$. By the uniqueness of primes lying over
$\mathfrak p'$ we see that $\mathfrak q' = \mathfrak q''$.  In other words
$\mathfrak q'$ defines a point of $\Spec(S_{\mathfrak q})$.

\medskip\noindent
In both cases we conclude that the map
$\Spec(S_{\mathfrak q}) \to \Spec(S_{\mathfrak p})$
is bijective. Clearly this means all the elements of $S - \mathfrak q$
are all invertible in $S_{\mathfrak p}$, in other words
$S_{\mathfrak p} = S_{\mathfrak q}$.
\end{proof}

\noindent
The following lemma is a generalization of going down for
flat ring maps.

\begin{lemma}
\label{lemma-going-down-flat-module}
Let $R \to S$ be a ring map. Let $N$ be a finite $S$-module flat over $R$.
Endow $\text{Supp}(N) \subset \Spec(S)$ with the induced topology.
Then generalizations lift along $\text{Supp}(N) \to \Spec(R)$.
\end{lemma}

\begin{proof}
The meaning of the statement is as follows. Let
$\mathfrak p \subset \mathfrak p' \subset R$ be primes. Let
$\mathfrak q' \subset S$ be a prime $\mathfrak q' \in \text{Supp}(N)$
Then there exists a prime $\mathfrak q \subset \mathfrak q'$,
$\mathfrak q \in \text{Supp}(N)$ lying over $\mathfrak p$.
As $N$ is flat over $R$ we see that $N_{\mathfrak q'}$ is flat
over $R_{\mathfrak p'}$, see Lemma \ref{lemma-flat-localization}.
As $N_{\mathfrak q'}$ is finite over $S_{\mathfrak q'}$
and not zero since $\mathfrak q' \in \text{Supp}(N)$ we see
that $N_{\mathfrak q'} \otimes_{S_{\mathfrak q'}} \kappa(\mathfrak q')$
is nonzero by Nakayama's Lemma \ref{lemma-NAK}.
Thus $N_{\mathfrak q'} \otimes_{R_{\mathfrak p'}} \kappa(\mathfrak p')$
is also not zero. We conclude from Lemma \ref{lemma-ff}
that $N_{\mathfrak q'} \otimes_{R_{\mathfrak p'}} \kappa(\mathfrak p)$
is nonzero. Let
$J \subset S_{\mathfrak q'} \otimes_{R_{\mathfrak p'}} \kappa(\mathfrak p)$
be the annihilator of the finite nonzero module
$N_{\mathfrak q'} \otimes_{R_{\mathfrak p'}} \kappa(\mathfrak p)$.
Since $J$ is a proper ideal we can choose a prime $\mathfrak q \subset S$
which corresponds to a prime of
$S_{\mathfrak q'} \otimes_{R_{\mathfrak p'}} \kappa(\mathfrak p)/J$.
This prime is in the support of $N$, lies over $\mathfrak p$, and
is contained in $\mathfrak q'$ as desired.
\end{proof}























